\documentclass[10pt,a4paper]{article}
\usepackage[T1]{fontenc}
\usepackage{lmodern}
\usepackage[margin=18mm,headheight=13pt,headsep=5mm,footskip=9mm]{geometry}
\usepackage{amsmath,graphicx,booktabs,caption,titlesec,fancyhdr,xcolor}
\definecolor{navy}{HTML}{173A55}
\usepackage[colorlinks=true,urlcolor=navy,linkcolor=navy]{hyperref}
\hypersetup{
  pdftitle={Assignment 2 - Interpolation and position measurement},
  pdfauthor={Paolo Rossi},
  pdfsubject={Multi Agent Systems for Smart Machining, ECIU Micro-Module}
}
\titleformat{\section}{\Large\sffamily\bfseries\color{navy}}{\thesection}{0.65em}{}
\titlespacing*{\section}{0pt}{8pt}{5pt}
\captionsetup{font=small,labelfont=bf,skip=4pt,hypcap=false}
\setlength{\parindent}{0pt}
\setlength{\parskip}{5pt}
\renewcommand{\arraystretch}{1.18}
\pagestyle{fancy}
\fancyhf{}
\fancyhead[L]{\small Assignment 2 \enspace$\cdot$\enspace Interpolation and position measurement}
\fancyhead[R]{\small Paolo Rossi -- 277177}
\fancyfoot[C]{\small\thepage}
\renewcommand{\headrulewidth}{0.3pt}
\graphicspath{{../plots/}}
\newcommand{\assignmenturl}{https://github.com/massma-students-material/assignment-02-paolo}
\newcommand{\repofile}[2]{\href{\assignmenturl/blob/79d9096c0ea39a8f55ba3f66c3150e4b9892b94f/#1}{#2}}
\newcommand{\authorprofile}{\href{https://github.com/PaoloRos}{Paolo Rossi}}
\newcommand{\reportfigure}[2]{%
  \par\medskip
  \noindent\begin{minipage}{\linewidth}
    \centering
    \includegraphics[width=\linewidth]{#1}
    \captionof{figure}{#2}
  \end{minipage}\par
}

\begin{document}
\thispagestyle{plain}
\begin{minipage}[c]{0.77\linewidth}
  {\Huge\sffamily\bfseries\color{navy} Assignment 2}\\[4pt]
  {\Large Interpolation and position measurement}
\end{minipage}\hfill
\begin{minipage}[c]{0.19\linewidth}
  \raggedleft\includegraphics[width=27mm]{unitn_logo.png}
\end{minipage}

Multi Agent Systems for Smart Machining \hfill \authorprofile{} -- 277177\\
ECIU Micro-Module, INSA Toulouse \hfill Academic Year 2026/27
\par\hrule

\section{Nominal feed and machining sequence}
The nominal feed along the milling contour is \textbf{955 mm/min (15.92 mm/s)}.
With a nominal cutter diameter of $4$ mm, three teeth, a cutting speed of $80$ m/min,
and a feed per tooth of $0.050$ mm/tooth, the spindle speed is
\[
  n=\frac{1000(80)}{\pi(4)}\simeq 6366\ \mathrm{rpm},
  \qquad v_f=0.050(3)n\simeq 955\ \mathrm{mm/min}.
\]
The milling cycle traverses the seven blocks in reverse order from the contour
listing in the CNC program. Both encoder traces were sampled every $4$ ms.
\par
\repofile{data/5ASM.MPF}{CNC program} \enspace$\cdot$\enspace
\repofile{README.md\#trace-data}{Trace inputs and signal definitions}

\begin{center}
\small
\captionof{table}{Phase intervals in seconds from the start of each recording.
  Points 1--2 use a tool length of $184.978$ mm and a plunge-end crossing of $-2.000$ mm.}
\label{tab:phases}\begin{tabular}{@{}lcc@{}}
\toprule
Phase & With material removal & Without material removal\\
\midrule
Positioning     & $0.000\div 28.830$  & $0.000\div 5.027$\\
Plunge          & $28.830\div 30.962$ & $5.027\div 5.172$\\
Milling         & $30.962\div 43.560$ & $5.172\div 14.848$\\
Retract         & $43.560\div 43.910$ & $14.848\div 15.199$\\
\bottomrule
\end{tabular}
\end{center}

Interpolate encoder Z crossings at the workpiece surface and the programmed depth
to mark the start and end of the plunge. The subsequent changes in commanded XY
targets mark the internal milling boundaries. The retract command ends milling.
Retract ends when encoder Z reaches within $0.02$ mm of the $50$ mm plane.
\par
\repofile{machining_intervals.py}{Interval code} \enspace$\cdot$\enspace
\repofile{outputs/machining_intervals.csv}{Phase and block timings}

\reportfigure{trajectory_3d.png}{Encoder trajectories and programmed block endpoints for both
  runs. Legend entries 1--3 identify plunge, milling, and retract using the timing
  convention in \hyperref[tab:phases]{Table~\ref*{tab:phases}}.}

\newpage
\section{X and Y feed components}
The $146.353$ mm contour takes \textbf{9.196 s nominally}, \textbf{12.598 s with
material removal}, and \textbf{9.676 s without}. Block 1 accounts for most of the
timing difference between the recorded runs. It takes $4.230$ s with material
removal and $1.316$ s without. The encoder shows a near-stop early in block 1 of
the material-removal run. The two runs have similar feed components at matching
contour positions in blocks 2--7.

\reportfigure{feed_velocity.png}{Nominal and encoder X and Y feed components by cumulative
  contour distance and elapsed milling time. For each recording, elapsed time
  starts at the $-2.000$ mm crossing with the $184.978$ mm tool length.}

The nominal XY velocity is the scalar feed $v_f$ times the unit tangent to the
executed contour. For a straight block with displacement $(\Delta x,\Delta y)$,
\[
  (v_x,v_y)=v_f\frac{(\Delta x,\Delta y)}{\sqrt{\Delta x^2+\Delta y^2}}.
\]
For the counterclockwise arc in block 2, the radius is $10$ mm and the arc begins
at $s=20$ mm. With $s$ measured from the start of the contour,
\[
  \phi=\frac{s-20\ \mathrm{mm}}{10\ \mathrm{mm}},
  \qquad (v_x,v_y)=v_f(\sin\phi,-\cos\phi).
\]
Encoder components come from the slopes of local quadratic fits over nine samples
spanning $32$ ms. Projecting each sample onto its assigned block gives the
cumulative distance $s$. The nominal time axis uses $t=s/v_f$.
\par
\repofile{feed_velocity.py}{Velocity code} \enspace$\cdot$\enspace
\repofile{outputs/feed_velocity.csv}{Velocity CSV} \enspace$\cdot$\enspace
\repofile{OUTCOMES.md\#point-2}{Nominal components for all seven blocks}

\newpage
\section{Machine-to-workpiece coordinates}
At the final retract sample with material removal ($44.432$ s), the converted
tool-tip position matches the programmed G55 position within $2.1~\mu\mathrm{m}$
on each axis.
\[
  \mathbf p_{\mathrm{meas}}=(-40.00029,-20.00017,50.00207)\ \mathrm{mm},
  \qquad \mathbf p_{\mathrm{prog}}=(-40,-20,50)\ \mathrm{mm}.
\]

Subtract the G55 offset and the axial tool length from the encoder position
$\mathbf p_M$ in machine coordinates. \hyperlink{slides}{Slide 8} gives
$\mathbf{G55}=(-279.872,-272.170,-374.244)$ mm and $D_1=184.970$ mm, so
\[
  \mathbf p_{\mathrm{G55}}
  =\mathbf p_M-\mathbf{G55}-(0,0,D_1)
  =\mathbf p_M+(279.872,272.170,189.274)\ \mathrm{mm}.
\]
Points 3--4 mark milling start at the first interpolated downward crossing of
$-1.990$ mm using the slide-8 tool length. That crossing occurs at $30.959$ s with
material removal and $5.170$ s without. Points 1--2 use the historical tool length
of $184.978$ mm and a $-2.000$ mm crossing.

\reportfigure{coordinates_vs_time.png}{Programmed targets and converted encoder positions
  over time. The phase markers and plunge details use the $-1.990$ mm threshold
  with the $184.970$ mm tool length.}

\reportfigure{trajectory_3d_origins.png}{Machine origin $O_M$ and G55 workpiece origin $O_P$ during the material-removal run. \repofile{plots/trajectory_3d_origins.png}{Full-size view}.}
\repofile{trace_data.py}{Conversion code} \enspace$\cdot$\enspace
\repofile{plot_traces.py}{Plotting code} \enspace$\cdot$\enspace
\repofile{outputs/workpiece_coordinates.csv}{Converted recordings}

\section{Signed XY contour error}
The largest component deviations are \textbf{$-28.86~\mu\mathrm{m}$ in X} with
material removal and \textbf{$-28.27~\mu\mathrm{m}$ in Y} without. Both occur at
the first sample of their respective blocks.

\begin{center}
\small
\captionof{table}{Largest absolute component error in each run, with the sign retained.}
\begin{tabular}{@{}lcccc@{}}
\toprule
Run & Axis and block & Error ($\mu$m) & Error (\% of block length) & Boundary flag\\
\midrule
With material removal    & X, block 2 & $-28.86$ & $-0.18373$ & First sample\\
Without material removal & Y, block 4 & $-28.27$ & $-0.28270$ & First sample\\
\bottomrule
\end{tabular}
\end{center}

\reportfigure{position_error.png}{Signed X and Y contour errors for both runs. Vertical lines
  mark the block boundaries and markers identify the largest absolute error in
  each axis, block, and run. Sample assignment uses the Point 3 intervals.}

Within each block, project the measured G55 XY sample onto the programmed line
segment or circular arc. For axis $i\in\{x,y\}$, subtract the corresponding nominal
coordinate $p_{i,\mathrm{proj}}$ and divide by the programmed block length $L_b$
when expressing the error as a percentage. Block 2 uses the arc length.
\[
  e_i=p_{i,\mathrm{enc}}-p_{i,\mathrm{proj}},
  \qquad e_{i,\%}=100\frac{e_i}{L_b}.
\]
The projection measures geometric deviation from the contour. Estimating
time-synchronized tracking error would require the interpolated controller setpoints.
\par
\repofile{position_error.py}{Error calculation code} \enspace$\cdot$\enspace
\repofile{outputs/position_error.csv}{Per-sample error CSV} \enspace$\cdot$\enspace
\repofile{OUTCOMES.md\#x-position-error}{X peak table} \enspace$\cdot$\enspace
\repofile{OUTCOMES.md\#y-position-error}{Y peak table}

\section{Engineering conclusions}
\begin{itemize}
  \item The two encoder traces show similar XY contour errors. Which run has the
    larger peak varies by axis and block. These recordings show no consistent
    increase in contour error during material removal.
  \item Small ripples appear on the arc and the two diagonal moves (blocks 2, 5,
    and 6). Across both runs their peak-to-peak errors range from $0.5$ to
    $1.1~\mu\mathrm{m}$ over the middle 80\% of those blocks.
  \item \textbf{19 of 28 per-axis peaks occur at block starts}, at the first encoder sample assigned to the new block. The errors generally return toward zero after these peaks.
  \item All absolute component errors are below \textbf{0.3\% of programmed
    block length} in the two runs.
\end{itemize}


\section{Comparison with colleagues}
The main difference is how we calculate the X and Y errors. My colleagues in
Group 2 in Toulouse compare the measured cutting and no-cutting paths directly.
For each segment and axis, they average the absolute differences between
corresponding positions on the two paths. I compare each run with the programmed
contour. For each encoder sample, I calculate the signed X and Y deviations
from the closest point on the programmed segment or arc, then compare these
deviations between the cutting and no-cutting runs.

\section*{AI usage}
\hyperlink{tasks}{\texttt{TODO.md}} specified the tasks. OpenAI's GPT-6 model through Codex generated
supporting material for the analysis and report. \hyperlink{outcomes}{\texttt{OUTCOMES.md}} recorded the
results. The author, \authorprofile, reviewed and approved all included material.

\section*{References}
{\small
\hypertarget{slides}{Anna Carla Araujo}, INSA Toulouse, \emph{Assignment 2: Interpolation and measuring
positions}, slides 3--4 and 8.
\repofile{assets/Assignment\%202\%20-\%20Slides.pdf}{Course slides}\par
Paolo Rossi, \emph{Assignment 2 analysis}.
\href{\assignmenturl}{Repository} \enspace$\cdot$\enspace
\hypertarget{tasks}{\repofile{TODO.md}{Task instructions}} \enspace$\cdot$\enspace \hypertarget{outcomes}{\repofile{OUTCOMES.md}{Analysis outcomes}}
}
\end{document}
